% from ega0-17-fr.tex
\oldpage[0\textsubscript{IV-1}]{48}
maximal $\mathfrak{m}'=\mathfrak{m}/aA$~; il est clair que
$\operatorname{rg}_k(\mathfrak{m}'/\mathfrak{m}'^2)=n-1$. En vertu de
(17.3.1.4), $\mathfrak{m}/a\mathfrak{m}$ est un $A'$-module de dimension
projective finie, donc il en est de même de
$\mathfrak{m}'=\mathfrak{m}/aA$, qui en est un facteur direct (17.3.1.3 et
17.2.1). L'hypothèse de récurrence entraîne par suite que $A'=A/aA$ est
régulier, ce qui prouve par (17.1.8) que $A$ est régulier.

\begin{corollary}[17.3.2]
\label{0.17.3.2-fr}
Si $A$ est un anneau local régulier, $A_{\mathfrak{p}}$ est régulier pour tout
idéal premier $\mathfrak{p}$ de $A$.
\end{corollary}

\begin{proof}
En effet, on a vu (17.2.10) que
$\operatorname{dim.coh}(A_{\mathfrak{p}})\leq
\operatorname{dim.coh}(A)$, donc la conclusion résulte aussitôt de (17.3.1).
\end{proof}

\begin{proposition}[17.3.3]
\label{0.17.3.3-fr}
Soient $A$, $B$ deux anneaux locaux noethériens,
$\varphi:A\longrightarrow B$ un homomorphisme local, $\mathfrak{m}$ (resp.
$\mathfrak{n}$) l'idéal maximal de $A$ (resp. $B$),
$k=A/\mathfrak{m}$ (resp. $k'=B/\mathfrak{n}$) le corps résiduel de $A$
(resp. $B$). On a donc un homomorphisme canonique de $k'$-espaces vectoriels
\begin{equation}
\label{0.17.3.3.1-fr}
\tag{17.3.3.1}
  \psi:(\mathfrak{m}/\mathfrak{m}^2)\otimes_k k'
  \longrightarrow\mathfrak{n}/\mathfrak{n}^2.
\end{equation}
\begin{enumerate}
  \item[\rm (i)] Si $B$ est régulier et si $B$ est un $A$-module plat, $A$
  est régulier.
  \item[\rm (ii)] Les conditions suivantes sont équivalentes~:
  \begin{enumerate}
    \item[\rm a)] $B$ est régulier et l'homomorphisme $\psi$ (17.3.3.1) est
    injectif.
    \item[\rm b)] $A$ et $B$ sont réguliers, et pour un système régulier de
    paramètres $(x_i)_{1\leq i\leq m}$ de $A$, les $\varphi(x_i)$ font partie
    d'un système régulier de paramètres de $B$ (auquel cas cette propriété a
    lieu pour \emph{tout} système régulier de paramètres de $A$).
    \item[\rm c)] $B$ et $B\otimes_A k$ sont réguliers, et $B$ est un
    $A$-module plat.
    \item[\rm d)] $A$ et $B\otimes_A k$ sont réguliers, et $B$ est un
    $A$-module plat.
    \item[\rm e)] $A$ et $B\otimes_A k$ sont réguliers, et l'on a
    \begin{equation}
    \label{0.17.3.3.2-fr}
    \tag{17.3.3.2}
      \dim(B)=\dim(A)+\dim(B\otimes_A k).
    \end{equation}
  \end{enumerate}
\end{enumerate}
\end{proposition}

\begin{proof}
\textnormal{(i)} On a
$\operatorname{dim.coh}(A)\leq\operatorname{dim.coh}(B)$ par (17.2.13), donc
il suffit d'appliquer (17.3.1).

\textnormal{(ii)} Lorsque $A$ est supposé régulier et que $(x_i)$ est un
système régulier de paramètres de $A$, dire que $B$ est un $A$-module
\emph{plat} équivaut, en vertu de (15.1.21) à dire que la suite des
$y_i=\varphi(x_i)$ est $B$-régulière (puisque $B\otimes_A k$ est un
$k$-module plat). D'autre part, comme $\dim(A)=m$, et que la suite $(y_i)$
engendre l'idéal $\mathfrak{m}B$, la relation (17.3.3.2), qui s'écrit aussi
$\dim(B/\mathfrak{m}B)=\dim(B)-m$, équivaut à dire que la suite
$(y_i)_{1\leq i\leq m}$ fait partie d'un système de paramètres de $B$
(16.3.7). On voit donc que les conditions b), d) et e) équivalent
respectivement aux suivantes~:
\begin{enumerate}
  \item[\rm b')] $A$ et $B$ sont réguliers et $(y_i)_{1\leq i\leq m}$ fait
  partie d'un système régulier de paramètres de $B$.
  \item[\rm d')] $A$ est régulier,
  $B/(\sum_{i=1}^m y_iB)$ est régulier et la suite $(y_i)$ est
  $B$-régulière.
  \item[\rm e')] $A$ est régulier,
  $B/(\sum_{i=1}^m y_iB)$ est régulier et la suite $(y_i)$ fait partie d'un
  système de paramètres de $B$.
\end{enumerate}
Or, b') et e') sont équivalentes en vertu de (17.1.7), et comme d') entraîne
e') (16.4.1) et est entraînée par b') (17.1.7), elle leur est équivalente. La
conjonction
